\begin{table*}[!th]
    \centering
    \small
    \caption{
        \textbf{Prompt used to rewrite \theorem questions.} {\color{blue} Blue text} denote instance-specific inputs. {\color{violet} Violet text} denotes the rewritten question and answer outputted by GPT-4.
        For each question, we provide the original question and answer, as well as the theorem name and theorem definition from the original TheoremQA dataset.
        The instruction prompts the model to rewrite the question with a different surface form without changing the reasoning steps or the final answer.
        We hand-write three examples to illustrate the rewriting process.
        We also allow the model to skip the question.
    }
    \label{tab:ex_rewrite}
    % \resizebox{0.98\linewidth}{!}{
    \begin{tabular}{>{\raggedright\arraybackslash\tt}p{0.98\textwidth}<{}}
        \toprule
            \vspace{-1em}
            Rewrite the following question such that the logical steps and final answer are still the same, only change the surface form of the question. Avoid using the words related to the theorem used in the question.
            You can do this by using real-world examples and applications to illustrate the concepts, in a way that is easier to understand for a layman.
            Try to ground the question in a real-world context such as finance and engineering problems, but be creative and feel free to use any domain you like! However, the question should be standalone and solvable.
            Rewrite the question in a json format, with the fields "question" and "answer".
            The "theorem" field indicates the theorem used in the question, but you should not use the words related to the theorem in the rewritten question.
            If you do not think that the question cannot be rewritten in layman's terms in a standalone fashion, then you can write a json object with the field "skip" set to True.
            \\\\
            \{
                "question": "In a group of 10 people, each of whom has one of 3 different eye colors, at least how many people must have the same eye color?",
                "answer": "4",
                "theorem": "pigeonhole principle",
                "theorem\_definition": "The Pigeonhole Principle is a fundamental concept in combinatorics, a branch of mathematics that deals with counting and arranging objects. \textit{[additional text omitted...]}"
            \}
            \\
            one demonstration was omitted...
            \\
            \{
                "question": "dy/dt = \textbackslash sqrt\{t\}, y(1) = 1. What is y(4)?",
                "answer": "5.667",
                "theorem": "integral rules",
                "theorem\_definition": "Integral rules in calculus are a set of techniques \textit{[additional text omitted...]} "
            \}
            \\
            \{
                "question": "You're tracking the growth of a plant from a seed. The rate at which the plant grows in height is equal to the square root of the number of days since you planted it. One day after the seed was first planted, the plant was 1 inch tall. How tall will the plant be after 4 days?",
                "answer": "5.667"
            \}
            \\\\
            \{
                % "question": "{\color{blue} \{Given image \backslash begin\{tabular\}\{\|llll\|\} \backslash hline 7 \& 1 \& 6 \& 0 \backslash 3 \& 3 \& 7 \& 6 \backslash 6 \& 6 \& 5 \& 7 \backslash \backslash hline \backslash end\{tabular\} , and the bit-depth of the image is 4. Suppose you want to use the thresholding technique to segment the image. What is the appropriate threshold value based on the histogram of the image? Follow the following rule when you do thresholding or grouping: pixel \$(i, j) \backslash in\$ Group A pixels if \$g(i, j) \backslash leq\$ current threshold \$\backslash mathrm\{T\}\$; pixel \$(i, j) \backslash in\$ Group B pixels otherwise, where \$g(i, j)\$ is the intensity value of pixel \$(i, j)\$.\} }",
                "question": "{\color{blue} \lstinline!Given image \begin{tabular}{|llll|} \hline 7 & 1 & 6 & 0 \\ 3 & 3 & 7 & 6 \\ 6 & 6 & 5 & 7 \\ \hline \end{tabular} , and the bit-depth of the image is 4. Suppose you want to use the thresholding technique to segment the image. What is the appropriate threshold value based on the histogram of the image? Follow the following rule when you do thresholding or grouping: pixel $(i, j) \in$ Group A pixels if $g(i, j) \leq$ current threshold $\mathrm{T}$; pixel $(i, j) \in$ Group B pixels otherwise, where $g(i, j)$ is the intensity value of pixel $(i, j)$.!}",
                "answer": "{\color{blue} 4}",
                "theorem": "{\color{blue} image contrast }",
                "theorem\_definition": "{\color{blue} Image contrast in signal processing refers to the difference in intensity or color between various elements or regions within an image. \textit{[additional text omitted...]} }"
            \}
            \\
            {\color{violet}\{
                "question": "Imagine you have a collection of 12 unique stickers, each with a number from 0 to 7 representing its rarity level. You're planning to divide these stickers into two groups based on their rarity to make it easier for collectors to understand. You decide to use a method where stickers with a rarity level at or below a certain number go into the 'Common' group, and those above this number go into the 'Rare' group. Given the distribution of stickers' rarity levels as follows: two stickers each of levels 7, 6, and 5; three stickers of level 3; and one sticker each of levels 1 and 0. What is the rarity level that should be used as the cutoff to divide the stickers into the 'Common' and 'Rare' groups, ensuring a balanced understanding of rarity?",
                "answer": "4"
            \}}
            \\
            % \noindent {\color{violet} \{ "question": "Imagine you have a collection of 12 unique stickers, each with a number from 0 to 7 representing its rarity level. You're planning to divide these stickers into two groups based on their rarity to make it easier for collectors to understand. You decide to use a method where stickers with a rarity level at or below a certain number go into the 'Common' group, and those above this number go into the 'Rare' group. Given the distribution of stickers' rarity levels as follows: two stickers each of levels 7, 6, and 5; three stickers of level 3; and one sticker each of levels 1 and 0. What is the rarity level that should be used as the cutoff to divide the stickers into the 'Common' and 'Rare' groups, ensuring a balanced understanding of rarity?", "answer": 4 \} }\\
        \bottomrule
    \end{tabular}
\end{table*}
